\documentclass{article}
\usepackage{amsmath, amssymb, amsthm}

\setlength{\parindent}{0pt}

\newtheorem*{claim}{Claim}

\begin{document}

\begin{claim}
    Let $X$ be a subspace of a vector space $V$, and let $\mathbf{v} \in V$, $\mathbf{v} \notin X$. Then 
    for every $\mathbf{x} \in X$, $\mathbf{x} + \mathbf{v} \notin X$.
\end{claim}

\begin{proof}
    Let $\mathbf{x} \in X$ be arbitrary. Suppose, for contradiction, that $\mathbf{x} + \mathbf{v} \in X$. 
    Since $X$ is closed under scalar multiplication, $(-1)\mathbf{x} \in X$, and 
    $(-1)\mathbf{x} = -\mathbf{x}$ by the vector space axioms. Furthermore, since $X$ is closed under 
    addition and $\mathbf{x} + \mathbf{v} \in X$, we get
    \[
    (\mathbf{x} + \mathbf{v}) + (-\mathbf{x}) \in X
    \]
    But using the vector space axioms, 
    \[
    \begin{aligned}
        (\mathbf{x} + \mathbf{v}) + (-\mathbf{x}) &= (\mathbf{v} + \mathbf{x}) + (-\mathbf{x}) \\
        &= \mathbf{v} + (\mathbf{x} + (-\mathbf{x})) \\
        &= \mathbf{v} + \mathbf{0} \\
        &= \mathbf{v}
    \end{aligned}
    \]
    Hence $\mathbf{v} \in X$, which is a contradiction. Since $\mathbf{x}$ was arbitrary, the claim holds 
    for all $\mathbf{x} \in X$.

    \medskip

    Hence $\mathbf{x} + \mathbf{v} \notin X$.
\end{proof}

\end{document}
